\subsection{NORMALIZE and GROUP}\label{subsec:knowledge-normalize-and-group}

Two deterministic stages run between the grounding stage and the canonical rule production: NORMALIZE and GROUP. The NORMALIZE stage standardizes entities within grounded findings, while GROUP groups normalized findings for canonicalization in the next stage. 

The NORMALIZE stage combines a curated synonym map -- consisting of hand-written mappings between common variants of biomedical terms -- with a biomedical entity linker ~\parencite{neumann2019scispacy, bodenreider2004umls} that resolves recognized entities to canonical names and UMLS Concept Unique Identifiers (CUIs). After entity normalization, a de-duplication pass merges findings that share the same six-field key: text-element id, normalized subject, normalized outcome, relation type, direction, and category. This process collapses each cluster into a single \texttt{NormalFinding}. The merged record preserves the source spans of each finding and carries the mean grounding score of the cluster. 

Direction and category are included in the deduplication key to prevent positive and negative claims, or claims with different directions, from being prematurely merged; this allows contradictions to be surfaced by a later RELATE stage. Keeping category information in the key prevents findings across different histopathological categories from being collapsed into a single category before GROUP has seen the data. A finding that misses any key field, such as a missing subject, outcome, or an unclear relation type, is not merged into anything, but is wrapped separately into its own \texttt{NormalFinding} object, so that nothing is dropped. 

The GROUP stage then puts \texttt{NormalFinding} records into buckets, per paper, with a key consisting of four fields: \texttt{subject\_entity}, \texttt{outcome\_entity}, \texttt{relation\_type}, and \texttt{category}. Each bucket produces one \texttt{FindingGroup}. At this stage, records with missing subjects, outcomes, or unclear relation types are dropped -- but they are logged in a rejection summary, rather than discarded silently. Since the process is run per paper, if two papers independently assert the same (subject, outcome, relation, category) 4-tuple, they will produce two distinct groups rather than one. Cross-paper merging is deferred to a separate corpus-level relation step over canonical rules (Subsection~\ref{subsec:knowledge-relate}), and is not performed in the GROUP stage.

The output of the GROUP stage is a list of \texttt{FindingGroup} records, each carrying its member \texttt{NormalFinding} records, which are passed down to the CANONICALIZE stage.

